
%%% fhc mu-like

\begin{table}[htb]
\centering
\begin{tabular}{| c | c | c |}
\hline
  {\bf Source of uncertainty}    & \multicolumn{1}{c|}{{\bf ${\delta}N_{SK}$/$N_{SK}$}}\\
                                 & \multicolumn{1}{c|}{(Asimov A)}\\

\hline
\hline
SKDet+FSI+SI                   & 4.11\%                        \\
SKDet only                   & 3.85\%                        \\
FSI+SI only                   & 1.44\%                        \\
\hline
\hline
  Flux                           & 3.61\%                        \\
\hline
\hline
  2p-2h (corr)                     & 3.53\%                        \\
  2p-2h-bar (corr)                  & 0.20\%                        \\
  NC other (uncorr)              & 0.78\%                        \\
  NC 1gamma (uncorr)             & 0.00\%                        \\
  XSec nue/numu (uncorr)                 & 0.01\%                        \\
\hline
  XSec Tot (corr)                  & 4.01\%                        \\
  XSec Tot                       & 4.09\%                        \\
\hline
\hline
  Flux+XSec (ND280 constrained)                      & 2.79\%                        \\
  Flux+XSec  (All)                    & 2.90\%                        \\
\hline
  Flux+XSec+SKDet+FSI+SI         & 5.09\%                        \\
  Flux+XSec+SKDet+FSI+SI (pre-fit)         & 12.2\%                        \\
\hline
\hline
  Oscillations                   & 0.03\%                        \\
\hline
\hline
  All                            & 5.17\%                        \\
  All (pre-fit) 		& 12.2\%				\\
\hline

\end{tabular}
\caption{
 Effect of 1$\sigma$ variations on the total number of events 1-ring muon-like events for neutrino mode
 given by all the systematic uncertainties, obtained by performing 10k toys MC.
The RMS / Mean of the distribution is assumed to be the $1\sigma$ uncertainty.
The systematic parameter values measured with the ND280 data 
shown in 
\cref{tab:banff_uncertainties_flux_fhc},
~\ref{tab:banff_uncertainties_flux_rhc}
~\ref{tab:banff_uncertainties_xsec}
and
~\ref{tab:banff_uncertainties_skdetfsi}
and the oscillation parameters of the Asimov data set A shown in \cref{tab:osc_params_joint_asimov}
are used.
Since the parameters $\delta_{CP}$, $\theta_{23}$ and $\Delta m^2_{32}$ ($\Delta m^2_{13}$) and mass hierarchy
will be measured,
only the uncertainty on the parameters 
$\sin^2 \theta_{13}$,
$\sin^2 \theta_{12}$
and 
$\Delta m^2_{21}$ 
are taken into account.
Pre-BANFF fit systematic uncertainties are used were stated.
The changes due to the Run 1-7b flux and covariance matrix updates are applied.
}                                                                                                                                                                                          
\label{tab:nsk_systematic_table_summary_joint_asimov_A_fhc_mulike_Run17b}
\end{table}




%%%%% fhc e-like


\begin{table}[htb]
\centering
\begin{tabular}{| c | c | c |}
\hline
  {\bf Source of uncertainty}    & \multicolumn{1}{c|}{{\bf ${\delta}N_{SK}$/$N_{SK}$}}\\
                                 & \multicolumn{1}{c|}{(Asimov A)}\\

\hline
\hline  
SKDet+FSI+SI                   & 3.48\%                        \\
SKDet only                   & 2.28\%                        \\
FSI+SI only                   & 2.63\%                        \\
\hline
\hline
  Flux                           & 3.67\%                        \\
\hline
\hline
  2p-2h (corr)                     & 3.90\%                        \\
  2p-2h bar (corr)                  & 0.05\%                        \\
  NC other (uncorr)              & 0.15\%                        \\
  NC 1gamma (uncorr)             & 1.47\%                        \\
  XSec nue/numu (uncorr)                 & 2.61\%                        \\
\hline
  XSec Tot (corr)                  & 4.26\%                        \\
  XSec Tot                       & 5.21\%                        \\
\hline
\hline
  Flux+XSec (ND280 constrained)                      & 2.90\%                        \\
  Flux+XSec  (All)                     & 4.17\%                        \\
\hline
  Flux+XSec+SKDet+FSI+SI         & 5.45\%                        \\
  Flux+XSec+SKDet+FSI+SI (pre-fit)         & 12.1\%                        \\
\hline
\hline
  Oscillations                   & 4.20\%                        \\
\hline
\hline
  All                            & 6.91\%                        \\
  All (pre-fit) 		& 12.6\%				\\
\hline
\hline

\end{tabular}
\caption{
 Effect of 1$\sigma$ variations on the total number of events 1-ring electron-like events for neutrino mode
 given by all the systematic uncertainties, obtained by performing 10k toys MC.
The RMS / Mean of the distribution is assumed to be the $1\sigma$ uncertainty.
The systematic parameter values measured with the ND280 data 
shown in 
\cref{tab:banff_uncertainties_flux_fhc},
~\ref{tab:banff_uncertainties_flux_rhc}
~\ref{tab:banff_uncertainties_xsec}
and
~\ref{tab:banff_uncertainties_skdetfsi}
and the oscillation parameters of the Asimov data set A shown in \cref{tab:osc_params_joint_asimov}
are used.
Since the parameters $\delta_{CP}$, $\theta_{23}$ and $\Delta m^2_{32}$ ($\Delta m^2_{13}$) and mass hierarchy
will be measured,
only the uncertainty on the parameters 
$\sin^2 \theta_{13}$,
$\sin^2 \theta_{12}$
and 
$\Delta m^2_{21}$ 
are taken into account.
Pre-BANFF fit systematic uncertainties are used were stated.
The changes due to the Run 1-7b flux and covariance matrix updates are applied.
}                                                                                                                                                                                          
\label{tab:nsk_systematic_table_summary_joint_asimov_A_fhc_elike_Run17b}
\end{table}



%%%%%% rhc mu-like


\begin{table}[htb]
\centering
\begin{tabular}{| c | c | c |}
\hline
  {\bf Source of uncertainty}    & \multicolumn{1}{c|}{{\bf ${\delta}N_{SK}$/$N_{SK}$}}\\
                                 & \multicolumn{1}{c|}{(Asimov A)}\\

\hline
\hline  
SKDet+FSI+SI                   & 3.94\%                        \\
SKDet only                   & 3.36 \%                        \\
FSI+SI only                   & 2.05 \%                        \\
\hline
  Flux                           & 3.80\%                        \\
\hline
\hline
  2p-2h (corr)                     & 2.96\%                        \\
  2p-2h bar (corr)                  & 1.83\%                        \\
  NC other (uncorr)              & 0.76\%                        \\
  NC 1gamma (uncorr)             & 0.00\%                        \\
  XSec nue/numu (uncorr)                 & 0.00\%                        \\
\hline
  XSec Tot (corr)                  & 4.13\%                        \\
  XSec Tot                       & 4.20\%                        \\
\hline
\hline
  Flux+XSec (ND280 constrained)                      & 3.27\%                        \\
  Flux+XSec  (All)                    & 3.36\%                        \\
\hline
  Flux+XSec+SKDet+FSI+SI         & 5.18\%                        \\
  Flux+XSec+SKDet+FSI+SI (pre-fit)         & 12.4\%                        \\
\hline
\hline
  Oscillations                   & 0.03\%                        \\
\hline
\hline
  All                            & 5.17\%                        \\
      All (pre-fit) 		& 12.5\%				\\
\hline
\hline
\end{tabular} 
\caption{
 Effect of 1$\sigma$ variations on the total number of events 1-ring muon-like events for antineutrino mode
 given by all the systematic uncertainties, obtained by performing 10k toys MC.
The RMS / Mean of the distribution is assumed to be the $1\sigma$ uncertainty.
The systematic parameter values measured with the ND280 data 
shown in 
\cref{tab:banff_uncertainties_flux_fhc},
~\ref{tab:banff_uncertainties_flux_rhc}
~\ref{tab:banff_uncertainties_xsec}
and
~\ref{tab:banff_uncertainties_skdetfsi}
and the oscillation parameters of the Asimov data set A shown in \cref{tab:osc_params_joint_asimov}
are used.
Since the parameters $\delta_{CP}$, $\theta_{23}$ and $\Delta m^2_{32}$ ($\Delta m^2_{13}$) and mass hierarchy
will be measured,
only the uncertainty on the parameters 
$\sin^2 \theta_{13}$,
$\sin^2 \theta_{12}$
and 
$\Delta m^2_{21}$ 
are taken into account.
Pre-BANFF fit systematic uncertainties are used were stated.
The changes due to the Run 1-7b flux and covariance matrix updates are applied.
}                                                                                                                                                                                          
\label{tab:nsk_systematic_table_summary_joint_asimov_A_rhc_mulike_Run17b}
\end{table}



%%%%%% rhc e-like

\begin{table}[htb]
\centering
\begin{tabular}{| c | c | c |}
\hline
  {\bf Source of uncertainty}    & \multicolumn{1}{c|}{{\bf ${\delta}N_{SK}$/$N_{SK}$}}\\
                                 & \multicolumn{1}{c|}{(Asimov A)}\\

\hline
\hline  
SKDet+FSI+SI                   & 3.95\%                        \\
SKDet only                   & 3.11\%                        \\
FSI+SI only                   & 2.43\%                        \\
\hline
  Flux                           & 3.84\%                        \\
\hline
\hline
  2p-2h (corr)                     & 3.04\%                        \\
  2p-2h bar (corr)                  & 2.36\%                        \\
  NC other (uncorr)              & 0.33\%                        \\
  NC 1gamma (uncorr)             & 2.95\%                        \\
  XSec nue/numu (uncorr)                 & 1.46\%                        \\
\hline
  XSec Tot (corr)                  & 4.46\%                        \\
  XSec Tot                       & 5.55\%                        \\
\hline
\hline
  Flux+XSec (ND280 constrained)                      & 3.20\%                        \\
  Flux+XSec                      & 4.60\%                        \\
\hline
  Flux+XSec+SKDet+FSI+SI         & 6.28\%                        \\
  Flux+XSec+SKDet+FSI+SI (pre-fit)         & 13.5\%                        \\
\hline
\hline
  Oscillations                   & 4.00\%                        \\
\hline
\hline
  All                            & 7.38\%                        \\
      All (pre-fit) 		& 14.1\%				\\
\hline
\hline
\end{tabular}
\caption{
 Effect of 1$\sigma$ variations on the total number of events 1-ring electron-like events for antineutrino mode
 given by all the systematic uncertainties, obtained by performing 10k toys MC.
The RMS / Mean of the distribution is assumed to be the $1\sigma$ uncertainty.
The systematic parameter values measured with the ND280 data 
shown in 
\cref{tab:banff_uncertainties_flux_fhc},
~\ref{tab:banff_uncertainties_flux_rhc}
~\ref{tab:banff_uncertainties_xsec}
and
~\ref{tab:banff_uncertainties_skdetfsi}
and the oscillation parameters of the Asimov data set A shown in \cref{tab:osc_params_joint_asimov}
are used.
Since the parameters $\delta_{CP}$, $\theta_{23}$ and $\Delta m^2_{32}$ ($\Delta m^2_{13}$) and mass hierarchy
will be measured,
only the uncertainty on the parameters 
$\sin^2 \theta_{13}$,
$\sin^2 \theta_{12}$
and 
$\Delta m^2_{21}$ 
are taken into account.
Pre-BANFF fit systematic uncertainties are used were stated.
The changes due to the Run 1-7b flux and covariance matrix updates are applied.
}                                                                                                                                                                                          
\label{tab:nsk_systematic_table_summary_joint_asimov_A_rhc_elike_Run17b}
\end{table}


















\begin{table}
\center
\begin{tabular}{| c || c | c || c | c |}
\hline
Source of uncertainty 
&$\mu$-like ${\delta} \Big( \frac{\# {\nu} \text{-mode}}{\# {\bar{\nu}} \text{-mode} } \Big) $ / $\Big \langle \frac{\# {\nu} \text{-mode}}{\# {\bar{\nu}} \text{-mode} } \Big \rangle$ 
&$e$-like ${\delta} \Big( \frac{\# {\nu} \text{-mode}}{\# {\bar{\nu}} \text{-mode} } \Big) $ / $\Big \langle \frac{\# {\nu} \text{-mode}}{\# {\bar{\nu}} \text{-mode} } \Big \rangle$ \\


\hline
\hline
SKDet 							& 0.7\%  	& 1.6\% \\ \hline
FSI+SI 	 						& 2.6\% 	& 3.5\% \\ \hline
SKDet+FSI+SI 						& 2.7\% 	& 3.9\% \\ \hline
\hline
Flux 								& 1.7\% 	& 1.7\% \\ \hline
XSec NC other (uncorr)  				& 0.0\% 	& 0.2\% \\ \hline
XSec NC $1\gamma$ (uncorr) 			& 0.0\%  	& 1.5\% \\ \hline
XSec $\nu_e$ / $\nu_{\mu}$ (uncorr)  	& 0.0\%  	& 3.1\% \\ \hline
XSec Tot (All)							& 1.7\%  	& 4.3\% \\ \hline
 \hline
Flux+XSec 						& 1.9\%  	& 4.1\% \\ \hline
 \hline
Flux+XSec+SDet+FSI+SI 				& 3.2\%  	& 5.9\% \\ \hline
 \hline
Oscillations 						& 0.0\% 	& 0.8\% \\ \hline
 \hline
All 								& 3.3\%  	& 5.8\% \\ \hline
\end{tabular}
\caption{
 Effect of 1$\sigma$ variations of the several source of systematic uncertainties 
 on the ratio between the
 total number of events in neutrino and antineutrino modes
 for 1-ring electron- and muon-like samples
 obtained by performing 10k toys MC.
 The systematic parameter values measured with the ND280 data 
shown in 
\cref{tab:banff_uncertainties_flux_fhc},
~\ref{tab:banff_uncertainties_flux_rhc}
~\ref{tab:banff_uncertainties_xsec}
and
~\ref{tab:banff_uncertainties_skdetfsi}
and the oscillation parameters of the Asimov data set A shown in \cref{tab:osc_params_joint_asimov}
are used.
Since the parameters $\delta_{CP}$, $\theta_{23}$ and $\Delta m^2_{32}$ ($\Delta m^2_{13}$) and mass hierarchy
will be measured,
only the uncertainty on the parameters 
$\sin^2 \theta_{13}$,
$\sin^2 \theta_{12}$
and 
$\Delta m^2_{21}$ 
are taken into account.
The changes due to the Run 1-7b flux and covariance matrix updates are applied.
}             
\label{tab:ratio_systematic_table_summary_asimov_A_Run17b}                                                                                                                                                                             
\end{table}
